% 第十三章：格林函数\index{格林函数}法

\chapter{格林函数\index{格林函数}法}

\begin{wulibj}[本章导读]
格林函数\index{格林函数}法是求解数学物理方程的一种强大而优美的方法。它源于物理学中的``叠加原理''：任何源产生的场可以看作点源产生的场的叠加。本章我们将学习格林函数\index{格林函数}的物理意义、数学推导以及它在求解各类边值问题中的应用。
\end{wulibj}

% ========== 第一节 ==========
\section{格林函数\index{格林函数}的概念与物理意义}

\subsection{从点电荷到一般电荷分布}

在静电学中，我们知道：
\begin{itemize}
    \item 点电荷 $q$ 在 $\boldsymbol{r}'$ 处产生的电势：$\phi(\boldsymbol{r}) = \frac{q}{4\pi\varepsilon_0|\boldsymbol{r} - \boldsymbol{r}'|}$
    \item 连续分布电荷产生的电势：$\phi(\boldsymbol{r}) = \frac{1}{4\pi\varepsilon_0}\int \frac{\rho(\boldsymbol{r}')}{|\boldsymbol{r} - \boldsymbol{r}'|}\,\mathrm{d}V'$
\end{itemize}

这就是格林函数\index{格林函数}思想的物理原型：先求出点源的响应，再通过积分（叠加）得到一般源分布的响应。



\subsection{$\delta$ 函数与格林函数\index{格林函数}}

点源在数学上用\textbf{$\delta$函数}表示。对于拉普拉斯方程，格林函数\index{格林函数}满足：
\begin{equation}
    \nabla^2 G(\boldsymbol{r}, \boldsymbol{r}') = -\delta(\boldsymbol{r} - \boldsymbol{r}')
\end{equation}

\begin{tishi}
$\delta$ 函数是一种广义函数，严格来说不是普通函数。它在积分号下有明确定义：
\begin{equation}
    \int_{-\infty}^{+\infty} f(x)\delta(x - a)\,\mathrm{d}x = f(a)
\end{equation}
\end{tishi}

% ========== 第二节 ==========
\section{拉普拉斯方程的格林函数\index{格林函数}}

\subsection{三维情形}

\begin{zhongyao}[三维拉普拉斯方程的基本解]
在三维无界空间中，格林函数为
\[
    G(\boldsymbol{r}, \boldsymbol{r}') = \frac{1}{4\pi|\boldsymbol{r} - \boldsymbol{r}'|}
\]
\end{zhongyao}

\begin{proof}
设 $\boldsymbol{R} = \boldsymbol{r} - \boldsymbol{r}'$，$R = |\boldsymbol{R}|$。

在 $R \neq 0$ 处，$\nabla^2\frac{1}{R} = 0$（验证：球坐标系下 $\nabla^2 f(R) = \frac{1}{R^2}\frac{\mathrm{d}}{\mathrm{d}R}\left(R^2\frac{\mathrm{d}f}{\mathrm{d}R}\right)$）。

在 $R = 0$ 处有奇点。利用高斯定理：
\begin{equation}
    \int_V \nabla^2\frac{1}{R}\,\mathrm{d}V = \oint_S \nabla\frac{1}{R} \cdot \mathrm{d}\boldsymbol{S} = \oint_S \left(-\frac{1}{R^2}\right)R^2\,\mathrm{d}\Omega = -4\pi
\end{equation}

所以
\begin{equation}
    \nabla^2\frac{1}{R} = -4\pi\delta(\boldsymbol{R}) = -4\pi\delta(\boldsymbol{r} - \boldsymbol{r}')
\end{equation}
\end{proof}

\subsection{二维情形}

\begin{zhongyao}[二维拉普拉斯方程的基本解]
在二维无界空间中，格林函数为
\[
    G(\boldsymbol{r}, \boldsymbol{r}') = \frac{1}{2\pi}\ln\frac{1}{|\boldsymbol{r} - \boldsymbol{r}'|} = -\frac{1}{2\pi}\ln|\boldsymbol{r} - \boldsymbol{r}'|
\]
\end{zhongyao}

\subsection{格林公式与格林函数\index{格林函数}表示}

\begin{dingli}{}{格林公式}
对于区域 $\Omega$ 内的函数 $u$ 和 $v$：
\begin{equation}
    \int_\Omega (u\nabla^2 v - v\nabla^2 u)\,\mathrm{d}V = \oint_{\partial\Omega}\left(u\frac{\partial v}{\partial n} - v\frac{\partial u}{\partial n}\right)\,\mathrm{d}S
\end{equation}
\end{dingli}

\begin{proof}
由高斯散度定理：
\begin{equation}
    \int_\Omega \nabla \cdot \boldsymbol{F}\,\mathrm{d}V = \oint_{\partial\Omega} \boldsymbol{F} \cdot \boldsymbol{n}\,\mathrm{d}S
\end{equation}
取 $\boldsymbol{F} = u\nabla v$，得
\begin{equation}
    \int_\Omega \nabla \cdot (u\nabla v)\,\mathrm{d}V = \oint_{\partial\Omega} u\frac{\partial v}{\partial n}\,\mathrm{d}S
\end{equation}
展开散度：$\nabla \cdot (u\nabla v) = \nabla u \cdot \nabla v + u\nabla^2 v$，故
\begin{equation}
    \int_\Omega u\nabla^2 v\,\mathrm{d}V = \oint_{\partial\Omega} u\frac{\partial v}{\partial n}\,\mathrm{d}S - \int_\Omega \nabla u \cdot \nabla v\,\mathrm{d}V
\end{equation}
交换 $u$ 和 $v$ 的位置，类似得到
\begin{equation}
    \int_\Omega v\nabla^2 u\,\mathrm{d}V = \oint_{\partial\Omega} v\frac{\partial u}{\partial n}\,\mathrm{d}S - \int_\Omega \nabla v \cdot \nabla u\,\mathrm{d}V
\end{equation}
两式相减即得格林公式。
\end{proof}
取 $v = G(\boldsymbol{r}, \boldsymbol{r}')$，利用格林公式可以证明：

\begin{dingli}{}{狄利克雷问题的解}
设 $u$ 在 $\Omega$ 内调和，在边界上给定 $u|_{\partial\Omega} = f$。则
\begin{equation}
    \boxed{u(\boldsymbol{r}') = -\frac{1}{4\pi}\oint_{\partial\Omega} f(\boldsymbol{r})\frac{\partial G}{\partial n}\,\mathrm{d}S}
\end{equation}
其中 $G$ 满足 $G|_{\partial\Omega} = 0$。
\end{dingli}

\begin{wulibj}[物理意义]
边界上的电势分布 $f$ 通过格林函数\index{格林函数}的影响，产生区域内的电势分布。格林函数\index{格林函数}起到了``传递''边界信息的作用。
\end{wulibj}

% ========== 第三节 ==========
\section{简单区域的格林函数\index{格林函数}}

\subsection{上半空间}

对于上半空间 $z > 0$，边界条件为 $G|_{z=0} = 0$。

使用\textbf{镜像法}：在 $(x', y', -z')$ 处放一个负电荷。

\begin{figure}[htbp]
\centering
\begin{tikzpicture}[scale=0.8]
    % 边界
    \draw[thick] (-3, 0) -- (3, 0);
    \fill[gray!30] (-3, -2) rectangle (3, 0);
    
    % 上半空间
    \fill[blue!10] (-3, 0) rectangle (3, 3);
    
    % 点源
    \fill[red] (0, 1.5) circle (3pt) node[right] {$(x', y', z')$};
    
    % 镜像
    \fill[blue] (0, -1.5) circle (3pt) node[right] {镜像电荷 $-1$};
    
    % 边界点
    \fill (2, 0) circle (2pt) node[above] {边界 $z=0$};
    
    \node at (0, 3.3) {上半空间};
    \node at (0, -1.8) {$z < 0$};
\end{tikzpicture}
\caption{上半空间的镜像法}
\label{fig:mirror-half-space}
\end{figure}

格林函数\index{格林函数}：
\begin{equation}
    \boxed{G(\boldsymbol{r}, \boldsymbol{r}') = \frac{1}{4\pi}\left[\frac{1}{|\boldsymbol{r} - \boldsymbol{r}'|} - \frac{1}{|\boldsymbol{r} - \boldsymbol{r}''|}\right]}
\end{equation}
其中 $\boldsymbol{r}'' = (x', y', -z')$ 为镜像点。

\subsection{球内区域}

对于球 $r < R$，边界条件为 $G|_{r=R} = 0$。

镜像电荷位置：$r_{\text{镜像}} = \frac{R^2}{r'}$，电荷量为 $-\frac{R}{r'}$。

格林函数\index{格林函数}：
\begin{equation}
    \boxed{G(\boldsymbol{r}, \boldsymbol{r}') = \frac{1}{4\pi}\left[\frac{1}{|\boldsymbol{r} - \boldsymbol{r}'|} - \frac{R/r'}{|\boldsymbol{r} - \frac{R^2}{r'^2}\boldsymbol{r}'|}\right]}
\end{equation}

\begin{lizi}{}{}
求上半空间 $z > 0$ 内拉普拉斯方程的解，边界条件 $u(x, y, 0) = f(x, y)$。

\begin{jie}
    由狄利克雷问题的解：
    \begin{equation}
        u(x', y', z') = -\frac{1}{4\pi}\int_{-\infty}^{+\infty}\int_{-\infty}^{+\infty} f(x, y)\frac{\partial G}{\partial z}\Big|_{z=0}\,\mathrm{d}x\,\mathrm{d}y
    \end{equation}
    
    计算 $\frac{\partial G}{\partial z}|_{z=0}$：
    \begin{equation}
        \frac{\partial G}{\partial z}\Big|_{z=0} = -\frac{z'}{2\pi[(x-x')^2 + (y-y')^2 + z'^2]^{3/2}}
    \end{equation}
    
    所以
    \begin{equation}
        \boxed{u(x', y', z') = \frac{z'}{2\pi}\int_{-\infty}^{+\infty}\int_{-\infty}^{+\infty} \frac{f(x, y)}{[(x-x')^2 + (y-y')^2 + z'^2]^{3/2}}\,\mathrm{d}x\,\mathrm{d}y}
    \end{equation}
    
    这就是\textbf{上半空间的泊松积分公式}。
\end{jie}
\end{lizi}

% ========== 第四节 ==========
\section{波动方程与热传导方程的格林函数\index{格林函数}}

\subsection{波动方程的格林函数\index{格林函数}}

对于三维波动方程的初值问题：
\begin{equation}
    \begin{cases}
        u_{tt} = a^2\nabla^2 u \\
        u(\boldsymbol{r}, 0) = 0, \quad u_t(\boldsymbol{r}, 0) = \psi(\boldsymbol{r})
    \end{cases}
\end{equation}

格林函数\index{格林函数}（推迟势）：
\begin{equation}
    G(\boldsymbol{r}, t; \boldsymbol{r}', t') = \frac{\delta(t - t' - |\boldsymbol{r} - \boldsymbol{r}'|/a)}{4\pi a^2|\boldsymbol{r} - \boldsymbol{r}'|}
\end{equation}

解为：
\begin{equation}
    u(\boldsymbol{r}, t) = \frac{1}{4\pi a^2}\int_{|\boldsymbol{r} - \boldsymbol{r}'| \leq at} \frac{\psi(\boldsymbol{r}')}{|\boldsymbol{r} - \boldsymbol{r}'|}\,\mathrm{d}V'
\end{equation}

\begin{wulibj}[因果性]
波动方程的格林函数\index{格林函数}只依赖于时间差 $t - t'$，体现了因果性——只有 $t' < t$ 的源才能影响 $t$ 时刻的场。
\end{wulibj}

\subsection{热传导方程的格林函数\index{格林函数}}

对于三维热传导方程的初值问题：
\begin{equation}
    \begin{cases}
        u_t = a^2\nabla^2 u \\
        u(\boldsymbol{r}, 0) = \varphi(\boldsymbol{r})
    \end{cases}
\end{equation}

\begin{zhongyao}[热传导方程的格林函数]
\[
    G(\boldsymbol{r}, t; \boldsymbol{r}', 0) = \frac{1}{(4\pi a^2 t)^{3/2}}\exp\left[-\frac{|\boldsymbol{r} - \boldsymbol{r}'|^2}{4a^2 t}\right]
\]
\end{zhongyao}

解为：
\begin{equation}
    u(\boldsymbol{r}, t) = \int_{\mathbb{R}^3} \varphi(\boldsymbol{r}')G(\boldsymbol{r}, t; \boldsymbol{r}', 0)\,\mathrm{d}V'
\end{equation}

这与傅里叶变换\index{傅里叶变换}法得到的泊松积分公式一致。

% ========== 本章小结 ==========
\section{本章小结}

\subsection{格林函数\index{格林函数}的定义}

格林函数\index{格林函数} $G(\boldsymbol{r}, \boldsymbol{r}')$ 是单位点源产生的场，满足：
\begin{equation}
    LG(\boldsymbol{r}, \boldsymbol{r}') = -\delta(\boldsymbol{r} - \boldsymbol{r}')
\end{equation}
其中 $L$ 是微分算子。

\subsection{基本结果}

\begin{table}[htbp]
\centering
\begin{tabular}{ll}
\toprule
\textbf{方程} & \textbf{格林函数\index{格林函数}（无界）} \\
\midrule
三维拉普拉斯方程 & $G = \frac{1}{4\pi|\boldsymbol{r} - \boldsymbol{r}'|}$ \\
二维拉普拉斯方程 & $G = -\frac{1}{2\pi}\ln|\boldsymbol{r} - \boldsymbol{r}'|$ \\
热传导方程 & $G = \frac{1}{(4\pi a^2 t)^{3/2}}\mathrm{e}^{-|\boldsymbol{r} - \boldsymbol{r}'|^2/(4a^2 t)}$ \\
\bottomrule
\end{tabular}
\end{table}

\subsection{求格林函数\index{格林函数}的方法}

\begin{enumerate}
    \item \textbf{镜像法}：适用于有简单几何对称性的区域
    \item \textbf{分离变量法\index{分离变量法}}：适用于规则区域
    \item \textbf{傅里叶变换\index{傅里叶变换}法}：适用于无界或半无界区域
\end{enumerate}

% ========== 习题 ==========
\section{习题}

\textbf{基础题}

\begin{enumerate}
    \item 验证 $G = \frac{1}{4\pi|\boldsymbol{r} - \boldsymbol{r}'|}$ 满足 $\nabla^2 G = -\delta(\boldsymbol{r} - \boldsymbol{r}')$。
    
    \item 证明二维拉普拉斯方程的基本解是 $G = -\frac{1}{2\pi}\ln|\boldsymbol{r} - \boldsymbol{r}'|$。
\end{enumerate}

\textbf{综合题}

\begin{enumerate}[resume]
    \item 用镜像法求球外区域 $r > R$ 的格林函数\index{格林函数}，边界条件 $G|_{r=R} = 0$。
    
    \item 求半无界区域 $x > 0$ 内拉普拉斯方程的格林函数\index{格林函数}，边界条件 $G|_{x=0} = 0$。
    
    \item 用格林函数\index{格林函数}法求解上半空间 $z > 0$ 的狄利克雷问题，边界条件 $u(x, y, 0) = 1$。
\end{enumerate}

\textbf{挑战题}

\begin{enumerate}[resume]
    \item 推导矩形区域 $0 < x < a$，$0 < y < b$ 内拉普拉斯方程的格林函数\index{格林函数}。
    
    \item 研究格林函数\index{格林函数}在量子力学散射理论中的应用（提示：与散射振幅的关系）。
\end{enumerate}
